\documentclass{report}
\usepackage{amsmath,amssymb}
\setlength{\parindent}{0mm}
\setlength{\parskip}{1em}
\begin{document}
\begin{center}
\rule{6in}{1pt} \
{\large Bart van Bloemen Waanders \\
{\bf Optimization under Uncertainty Constrained by Multiscale Phenomena}}

Sandia National Laboratories \\ PO Box 5800 \\ MS 0370 \\ Albuquerque \\ NM 87112
\\
{\tt bartv@sandia.gov}\\
Timothy Wildey\\
D. Thomas Seidl\end{center}

Many computational challenges in engineering and science exhibit
multiscale and multiphysics phenomena so that large scale optimization
and uncertainty quantification are computationally very expensive or
completely intractable. Furthermore, uncertainties arising at each scale
must be propagated through appropriate mapping mechanisms between scales.
In this work, we investigate multiscale formulations to improve on the
mapping mechanisms, apply computationally efficient optimization methods
to address the significant increase in degrees of freedom, and explore
mechanisms to handle uncertainties. We present a computational framework
based on multiscale mortar non-overlapping domain decomposition, coupled
to Newton-Krylov based optimization and endowed with concepts from
stochastic optimization. The multiscale mortar technique provides a
convenient and mathematically elegant approach for coupling different
numerical methods through physically meaningful interface conditions.
Neighboring sub-grids are not required to match along the interfaces.
This approach has been used to simulate multiphase flow through porous
media, compressible fluid dynamics, multiphysics, computational
mechanics, and geomechanics. The large scale optimization strategies we
employ are capable of addressing model uncertainties through risk
measures, which dictate performance criteria on the handling of the
uncertainties in the optimization process. We demonstrate our
methodologies on linear elasticity but note that our computational
framework is sufficiently general to handle a wide range of applications.


\end{document}
